% !TEX program = xelatex
% Compile with: latexmk -xelatex -interaction=nonstopmode technical_roadmap.tex
\documentclass[11pt,a4paper]{article}

\usepackage[a4paper,margin=2.4cm]{geometry}
\usepackage{xeCJK}
\usepackage{amsmath,amssymb}
\usepackage{booktabs}
\usepackage{longtable}
\usepackage{array}
\usepackage{tabularx}
\usepackage[hyphens]{url}
\usepackage{hyperref}
\usepackage{xcolor}
\usepackage{listings}
\usepackage{enumitem}
\usepackage{tikz}
\usepackage{caption}
\usepackage{titlesec}
\usepackage{fancyhdr}

\usetikzlibrary{positioning,arrows.meta,shapes.geometric,fit,backgrounds,calc}

% --- 中文字体（Windows 默认有 Microsoft YaHei / SimSun；Linux 可换 Noto） ---
\IfFontExistsTF{Microsoft YaHei}{%
  \setCJKmainfont[BoldFont={SimHei},ItalicFont={KaiTi}]{SimSun}
  \setCJKsansfont{Microsoft YaHei}
  \setCJKmonofont{NSimSun}
}{%
  \setCJKmainfont[BoldFont={Noto Sans CJK SC Bold}]{Noto Serif CJK SC}
  \setCJKsansfont{Noto Sans CJK SC}
  \setCJKmonofont{Noto Sans Mono CJK SC}
}

% --- Layout ---
\setlength{\parindent}{0pt}
\setlength{\parskip}{0.5em}
\linespread{1.18}

\titleformat{\section}{\Large\bfseries\sffamily}{\thesection.}{0.6em}{}
\titleformat{\subsection}{\large\bfseries\sffamily}{\thesubsection}{0.5em}{}
\titleformat{\subsubsection}{\normalsize\bfseries\sffamily}{\thesubsubsection}{0.4em}{}

\pagestyle{fancy}
\fancyhf{}
\fancyhead[L]{\small contract-review-openclaw}
\fancyhead[R]{\small 技术路线 v1.1}
\fancyfoot[C]{\thepage}
\renewcommand{\headrulewidth}{0.3pt}

\hypersetup{%
  colorlinks=true,
  linkcolor={black!75},
  urlcolor={blue!55!black},
  citecolor={black!75},
  pdftitle={contract-review-openclaw 技术路线},
  pdfauthor={liangch97},
}

% --- listings 中文友好 ---
\definecolor{codebg}{RGB}{248,248,250}
\definecolor{codecmt}{RGB}{106,153,85}
\definecolor{codekw}{RGB}{0,90,156}
\definecolor{codestr}{RGB}{163,21,21}

\lstdefinestyle{plain}{%
  basicstyle=\ttfamily\small,
  backgroundcolor=\color{codebg},
  frame=single,
  framesep=4pt,
  rulecolor=\color{black!18},
  breaklines=true,
  breakatwhitespace=false,
  columns=fullflexible,
  keepspaces=true,
  showstringspaces=false,
  xleftmargin=6pt,
  xrightmargin=6pt,
  literate={─}{{-}}1 {└}{{|-}}1 {├}{{|-}}1 {│}{{|}}1 {一}{{一}}1 {二}{{二}}1 {三}{{三}}1
}
\lstdefinestyle{py}{style=plain,language=Python,
  keywordstyle=\color{codekw}\bfseries,
  commentstyle=\color{codecmt}\itshape,
  stringstyle=\color{codestr}}
\lstdefinestyle{json}{style=plain,
  keywordstyle=\color{codekw}\bfseries,
  commentstyle=\color{codecmt}\itshape,
  stringstyle=\color{codestr}}
\lstdefinestyle{shell}{style=plain,
  keywordstyle=\color{codekw}\bfseries,
  commentstyle=\color{codecmt}\itshape}
\lstset{style=plain}

\title{\vspace{-1.4cm}\textbf{contract-review-openclaw 技术路线}\\[2pt]
  \large 形式化审核流水线 \quad 范式合同对照 skill 嵌入 \quad 双门槛范本匹配器}
\author{liangch97}
\date{2026-05-07 \quad v1.1}

\begin{document}
\maketitle
\thispagestyle{fancy}

\begin{abstract}\noindent
本文档描述 \texttt{contract-review-openclaw-portable} 的技术路线，涵盖从合同输入
到行政 PDF 输出的全流程；说明本次（v1.1）把范式合同对照能力拆为独立 OpenClaw
skill 的设计动机与具体方案；给出范本匹配器从「单一形状相似度」升级为
「形状 $+$ 锚点」双门槛后的判定逻辑、阈值取值和回归验证结果。文末附上后续
迭代规划。
\end{abstract}

\tableofcontents
\vspace{0.5em}

% =============================================================
\section{背景与目标}

中山大学横向科研合同的形式化审核长期存在两类问题：

\begin{enumerate}[leftmargin=1.4em,itemsep=2pt]
  \item \textbf{形式化检查的覆盖度}：行政老师手工核对甲方信息、关键条款、知识
    产权归属、保密期限等，工作量大且容易疏漏。
  \item \textbf{学校范本符合性}：中山大学发布了多份合同范本（技术开发委托、
    技术服务、技术咨询、专利技术转让、共同申请专利等），实际签订的合同与范本
    的偏离程度需要逐条比对。
\end{enumerate}

\texttt{contract-review-openclaw-portable} 是面向 OpenClaw 平台的本地包，目标
是让\textbf{行政老师只在 IM 中发一句话 + 一个合同附件}就拿到一份可直接转发的
\emph{行政版 PDF 报告}，包含：

\begin{itemize}[leftmargin=1.4em,itemsep=1pt]
  \item 形式化审核结论（基本可审 / 建议修改 / 存在重大问题）。
  \item 关键问题清单与所对应合同条款 / 证据。
  \item 甲方企业核验结果（企查查 / 人工材料 / 未完成核验）。
  \item 知识产权归属与背景知识产权使用安排。
  \item \textbf{学校范本对照}：是否基于学校范本，偏离了哪些条款。
\end{itemize}

% =============================================================
\section{整体架构}

\subsection{从一句话请求到行政 PDF}

图~\ref{fig:overall} 给出端到端流水线。所有模块均位于本地 OpenClaw / WSL，
无外部 LLM 调用，企查查访问受用户控制（扫码登录或本机会话复用）。

\begin{figure}[h]
\centering
\begin{tikzpicture}[
  node distance=10pt and 14pt,
  every node/.style={font=\small},
  box/.style={draw,rounded corners=2pt,thick,minimum width=2.6cm,
              minimum height=0.85cm,align=center,fill=white},
  data/.style={draw,thick,fill=black!4,minimum width=2.6cm,
               minimum height=0.85cm,align=center},
  skill/.style={draw,thick,fill=blue!8,minimum width=3.4cm,
                minimum height=0.85cm,align=center},
  outbox/.style={draw,thick,fill=green!8,minimum width=2.8cm,
                 minimum height=0.85cm,align=center},
  arr/.style={-{Latex[length=2.6mm]},thick}
]
\node[skill] (s1)  {OpenClaw skill\\\texttt{contract-formal-review-flow}};
\node[box,below=of s1] (full)  {\texttt{full\_qcc\_review.py}\\包装器};
\node[box,below=of full] (run) {\texttt{run\_contract\_review.py}\\本地命令行};

\node[box,right=2.4cm of full] (qcc) {QCC 扫码 /\\本机会话复用};
\node[box,below=of run] (load)  {\texttt{contract\_loader.py}};
\node[box,below=of load] (ext)  {\texttt{contract\_extractor.py}};
\node[box,below=of ext]  (tm)   {\texttt{template\_matcher.py}\\双门槛};
\node[box,below=of tm]   (rule) {\texttt{contract\_rule\_checker.py}};
\node[box,below=of rule] (rep)  {\texttt{contract\_report\_writer.py}};

\node[data,right=2.4cm of ext]  (jdata)  {\texttt{contract\_findings\_*.json}\\含 \texttt{template\_match}};
\node[skill,right=2.4cm of tm]  (s2)     {OpenClaw skill\\\texttt{contract-template-compliance}};
\node[outbox,right=2.4cm of rule]  (pdf)    {行政版 PDF};
\node[outbox,right=2.4cm of rep]   (mdxlsx) {Markdown + Excel};

\draw[arr] (s1) -- (full);
\draw[arr] (full) -- (run);
\draw[arr] (full) -- (qcc);
\draw[arr] (run) -- (load);
\draw[arr] (load) -- (ext);
\draw[arr] (ext) -- (tm);
\draw[arr] (tm) -- (rule);
\draw[arr] (rule) -- (rep);

\draw[arr] (rep) -- (mdxlsx);
\draw[arr] (rule) -- (pdf);
\draw[arr,dashed] (ext) -- (jdata);
\draw[arr,dashed] (tm)  -- (jdata);
\draw[arr,dashed] (jdata) -- (s2);

\end{tikzpicture}
\caption{端到端流水线：两个并列 skill、一条管线、一份 PDF。}
\label{fig:overall}
\end{figure}

\subsection{两个 skill 的分工}

\textbf{设计原则}：每个 OpenClaw skill 单一职责、$\le$ 5 KB description；当一个
skill 的 \texttt{SKILL.md} 体量超过 15 KB 或 covering 5+ 异构工作流时，模型
attention 会被稀释、出现遗忘和混淆。本项目在 v1.1 迭代前正是因为把范本对照
说明塞进主 skill，导致 \texttt{SKILL.md} 涨到 22 KB / 277 行，并发触发了
「忘记 IM 模式回复结构」「忘记调用 template-match」等 skill 把握能力下降的现象。

\begin{longtable}{p{4.6cm} p{3.3cm} p{6.8cm}}
\toprule
\textbf{Skill} & \textbf{规模} & \textbf{职责} \\ \midrule
\endhead
\texttt{contract-formal-review-flow} & 11.4 KB / 214 行
& 形式化审核默认入口；执行规则、IM 模式、最终回复结构、知识产权摘要规则；
  \emph{不}重述范本对照逻辑。 \\
\texttt{contract-template-compliance} & 7.2 KB / 99 行
& 范式合同对照专家：读取 \texttt{template\_match} 字段并产出范本符合性结论；
  \emph{不}独立运行任何脚本。 \\
\bottomrule
\end{longtable}

% =============================================================
\section{核心流水线}

\subsection{合同读取（\texttt{contract\_loader.py}）}

支持 \texttt{.docx} / \texttt{.pdf} / \texttt{.txt}，优先 \texttt{.docx}。
\texttt{.docx} 路径会读取：

\begin{itemize}[leftmargin=1.4em,itemsep=1pt]
  \item 正文段落 \texttt{paragraphs}
  \item 表格内容（含合并单元格按 \texttt{|} 拼接）
  \item Word 批注 \texttt{word/comments.xml}
  \item 高亮样式 \texttt{w:highlight}
\end{itemize}

输出 \texttt{ordered\_blocks}（按文档顺序的段落 / 表格混合块）+ \texttt{full\_text}。

\subsection{字段抽取（\texttt{contract\_extractor.py}）}

v1.1 在抽取层做了三件事：

\begin{enumerate}[leftmargin=1.4em,itemsep=2pt]
  \item \textbf{放宽标题正则}。原 \texttt{HEADING\_RE} 只匹配 \texttt{第 X 条} /
    阿拉伯数字 /\texttt{附件}；新增对 \texttt{一、} /\texttt{（一）} /
    \texttt{第一部分} 等编号的支持，覆盖 SYSU 范本的「使用说明」编号。
  \item \textbf{党事方角色别名}。原代码只识别 \texttt{甲方} /\texttt{乙方}；
    范本里大量使用 \emph{委托方} /\emph{受托方} /\emph{转让方} /\emph{受让方} /
    \emph{需求方} /\emph{服务方} /\emph{合作方} /\emph{供方}，因此引入
    \texttt{\_PARTY\_ROLE\_ALIASES} 把它们映射到 \texttt{甲} /\texttt{乙}
    抽取槽位。
  \item \textbf{接入范本匹配器}。\texttt{\_detect\_template\_match\_safe} 包
    装异常并返回 \texttt{\{matched: False, error: ...\}}，把
    \texttt{template\_match} 写入 \texttt{extracted}。
\end{enumerate}

\subsection{规则引擎（\texttt{contract\_rule\_checker.py}）}

规则文件 \texttt{references/rules/horizontal\_contract\_formal\_rules.yaml}，
按严重度分三级：

\begin{itemize}[leftmargin=1.4em,itemsep=1pt]
  \item \texttt{enforce}：自动可判定的硬性风险（如金额未填、知识产权归属缺失）。
  \item \texttt{review}：需要人工确认的事项（如保密期限是否合规）。
  \item \texttt{not\_checked}：仅作记录，\textbf{不}计入风险数。
\end{itemize}

\subsection{报告写出（\texttt{contract\_report\_writer.py}）}

输出三件套：\texttt{contract\_review.md}（行政摘要）、\texttt{contract\_findings\_<hash>.json}
（结构化结果，含 \texttt{template\_match}、\texttt{summary\_stats}、\texttt{findings}）、
\texttt{合同形式化审核总表.xlsx}。

经 \texttt{full\_qcc\_review.py} 触发时，额外渲染并通过飞书推送
\texttt{合同审核报告\_行政版.pdf}（即行政老师拿到的最终交付物）。

% =============================================================
\section{范本匹配器：从单门槛到双门槛}

\subsection{问题}

v1.0 的 \texttt{template\_matcher.py} 用 \texttt{difflib.SequenceMatcher} 直接
对「合同正文 vs 范本正文（去波动 token 后）」算相似度，超过
$\theta_{\text{shape}}=0.40$ 即判定为「基于学校范本」。

\textbf{失败案例}（反例）：行业起草的技术服务合同，乙方为中山大学。该合同与
\texttt{中山大学技术服务合同.docx} 共享大量套话（甲乙方、知识产权、保密、违约
责任、争议解决），但其余条款完全是行业自己的。结果：

\begin{itemize}[leftmargin=1.4em,itemsep=1pt]
  \item 实测相似度 $\approx 0.56$，被误判为「基于学校范本」。
  \item 行政 PDF 的「三、学校范本对照」一节里凭空冒出 28 个 \texttt{removed}、
    32 个 \texttt{added}（来自范本前置「使用说明」段被切成条款的噪声）。
\end{itemize}

\subsection{双门槛设计}

新版判定要求\textbf{形状信号与内容信号同时通过}：

\begin{equation}\label{eq:gate}
\text{matched} \;=\;
\bigl(s_\text{best} \ge \theta_\text{shape}\bigr)
\;\wedge\;
\bigl(\#\text{anchor}_\text{contract} \ge k_\text{anchor}\bigr)
\end{equation}

其中：

\begin{itemize}[leftmargin=1.4em,itemsep=1pt]
  \item $s_\text{best}=\max_T \text{Ratio}\bigl(\sigma(C),\sigma(T)\bigr)$，
    $C$ 是合同文本、$T$ 是范本文本、$\sigma$ 是
    \texttt{\_strip\_volatile}（清掉日期 / 金额 / 占位符等易变 token）。
  \item $\theta_\text{shape}=0.65$（v1.0 是 0.40；v1.1 抬高 25 个百分点，
    在反例 0.56 与正例 1.00 之间稳定区分）。
  \item \texttt{anchor} 集为\textbf{学校范本独有}的短语集合（见
    表~\ref{tab:anchors}），$k_\text{anchor}=1$。
\end{itemize}

\begin{table}[h]
\centering
\caption{SYSU 范本独有锚点短语（v1.1）}\label{tab:anchors}
\begin{tabular}{p{8.0cm} p{6.0cm}}
\toprule
\textbf{保留（独有）} & \textbf{已剔除（中性 / 误导）} \\ \midrule
\texttt{【科学研究院】} & \texttt{中山大学} \\
\texttt{【科研管理部门】} & \texttt{中大} \\
\texttt{学校合同管理信息系统} & \texttt{示范文本} \\
\texttt{合同管理信息系统} & \\
\texttt{中山大学·深圳} / \texttt{中山大学•深圳} / \texttt{中山大学∙深圳} & \\
\texttt{本合同由【科学研究院】} & \\
\texttt{本合同由【科研管理部门】} & \\
\bottomrule
\end{tabular}
\end{table}

\textbf{为何剔除 \texttt{中山大学} / \texttt{示范文本} / \texttt{中大}}？
中大经常作为乙方出现在行业自拟合同中（这只能说明合同的乙方是中大，并不意味着
合同采用了\emph{学校范本}）；\texttt{示范文本} 还会出现在国家科技部范本里。
这三个短语没有判别力，必须移出锚点集合。

\subsection{使用说明剥离（\texttt{\_strip\_preamble}）}

SYSU 范本 \texttt{.docx} 起始处有一段「使用说明」，编号为
\texttt{一、二、三、…七、}，解释如何填写范本。这些不是真正的合同条款，会在
逐条 diff 时产生大量 \texttt{removed} / \texttt{added} 噪声。

实现策略：当文本中存在 \texttt{第 N 条} 时，丢弃第一处 \texttt{第 N 条} 之前
的所有行；否则保留原文（避免误删纯条目式合同）。

\subsection{决策结果与回归验证}

实测两个样例（通过 \texttt{run\_contract\_review.py} 全流水线）：

\begin{longtable}{p{4.0cm} p{1.2cm} p{1.4cm} p{1.4cm} p{6.4cm}}
\toprule
\textbf{案例} & \textbf{matched} & \textbf{相似度} & \textbf{anchor 命中}
& \textbf{结论} \\ \midrule
\endhead
反例 \texttt{PCDW\_neg.docx}（行业自拟、乙方为中大）
& False & 0.41 & 0
& \texttt{non\_template\_contract: similarity\_0.41<0.65 and anchors\_0<1} \\
正例 \texttt{zhuanli\_pos.docx}（基于专利转让范本）
& True & 0.95 & 5
& 匹配《中山大学专利技术转让合同.docx》；未改动 38 / 修改 9 / 改写 5 /
  新增 2 / 删除 7 \\
\bottomrule
\end{longtable}

% =============================================================
\section{Skill 拆分细节}

\subsection{\texttt{contract-formal-review-flow}（v1.1 精简版）}

主入口 skill。\textbf{保留}：

\begin{itemize}[leftmargin=1.4em,itemsep=1pt]
  \item 强制执行规则（每次必须重跑脚本、不输出文本摘要、单条飞书回复等）。
  \item 禁用响应表（含「未收到附件」「请告诉我您的称呼」等典型错误）。
  \item IM 模式默认命令 \texttt{python3 /root/full\_qcc\_review.py "<path>"}。
  \item 最终回复 7 节结构（含「学校范本对照」一节，但只占一句话）。
  \item 知识产权摘要规则。
\end{itemize}

\textbf{外移}至 \texttt{references/handoff/qcc\_and\_arkclaw.md}：QCC 二维码登录
demo、ArkClaw 连接器输出清单、本地浏览器会话复用细节。

\textbf{YAML \texttt{description}}：双引号字符串，规避 \texttt{:\ } 触发的
mapping 歧义。

\subsection{\texttt{contract-template-compliance}（v1.1 新建）}

\textbf{触发词}：范本 / 范式 / 示范文本 / 学校范本 / 中大范本 / 范本对照 /
模板对照 / 条款偏离 / template compliance / clause deviation。

\textbf{决策规则}（必须三选一）：

\begin{enumerate}[leftmargin=1.4em,itemsep=2pt]
  \item $\texttt{matched}=\text{False}$：「本合同不是基于中山大学合同范本起草，
    按通用横向合同审核要点审查」。
  \item $\texttt{matched}=\text{True} \wedge \texttt{modified\_count}=0$：
    「与学校范本《\texttt{<name>}》一致，无条款偏离」。
  \item $\texttt{matched}=\text{True} \wedge \texttt{modified\_count}>0$：
    给出范本名 + 偏离统计 + 重点偏离条款表（按 \texttt{diff\_ratio} 升序最多 6 条）。
\end{enumerate}

\textbf{禁止行为}：

\begin{itemize}[leftmargin=1.4em,itemsep=1pt]
  \item 在 \texttt{matched=False} 时列「范本对照表」。
  \item 仅凭乙方是「中山大学」就声称采用了学校范本。
  \item 把 \texttt{template\_match.reason} 字段（含英文）原样贴给老师。
  \item 直接显示英文 status（应映射为「未改动 / 有修改 / 重大改写 /
    合同新增 / 范本中存在但合同未见」）。
\end{itemize}

\subsection{\texttt{template\_match} 数据契约}

\begin{lstlisting}[style=json,caption={\texttt{template\_match} 字段（v1.1）}]
{
  "matched": true,
  "template_name": "中山大学专利技术转让合同.docx",
  "similarity": 0.947,
  "anchor_hits": 5,
  "anchor_min": 1,
  "match_threshold": 0.65,
  "templates_considered": [...],
  "reason": null,
  "clauses": [
    {
      "id": "第三条",
      "title": "...",
      "status": "modified",
      "diff_ratio": 0.72,
      "template_excerpt": "...",
      "contract_excerpt": "..."
    }
  ],
  "counts": {"unchanged": 38, "modified": 9, "rewritten": 5,
             "added": 2, "removed": 7},
  "modified_count": 23,
  "summary": "匹配范本：…（相似度 95%，命中范本特征 5 项）；…"
}
\end{lstlisting}

% =============================================================
\section{部署与运维}

\subsection{安装}

\begin{lstlisting}[style=shell]
# 1. 克隆仓库
git clone https://github.com/liangch97/verifyAgent.git
cd verifyAgent

# 2. Python 依赖
pip install -r requirements.txt

# 3. 安装两个 skill 到 ~/.openclaw/skills/
python install_into_openclaw.py
\end{lstlisting}

\subsection{回归}

\begin{lstlisting}[style=shell]
# 范本匹配器单元
python3 _diag_template_match.py
python3 _diag_anchors.py
python3 _diag_clauses.py

# Skill 元数据 + 匹配器联合
python3 _smoketest_skills.py

# 端到端流水线
python3 _e2e_template.py
\end{lstlisting}

\subsection{回滚}

如需回滚到 v1.0 的整段 SKILL.md：

\begin{lstlisting}[style=shell]
cp ~/.openclaw/skills/contract-formal-review-flow/SKILL.md.preTemplateSplit.bak \
   ~/.openclaw/skills/contract-formal-review-flow/SKILL.md
rm -rf ~/.openclaw/skills/contract-template-compliance
\end{lstlisting}

% =============================================================
\section{后续迭代规划}

\begin{enumerate}[leftmargin=1.4em,itemsep=2pt]
  \item \textbf{语义化范本匹配}：在 \texttt{template\_matcher.py} 中加入条款级
    embedding（如 BGE-small-zh）做「语义相似度」与现有形状相似度的融合，进一
    步降低边缘案例（相似度 0.55$\sim$0.65）的误判风险。
  \item \textbf{条款级建议生成}：把 \texttt{clauses[].status=rewritten} 的条款
    交给规则引擎做二次审查，生成针对性的修改建议。
  \item \textbf{多范本组合}：支持「主范本 + 附加条款范本」的组合识别（如
    专利技术转让 + 共同申请专利的复合合同）。
  \item \textbf{打分一致性测试}：构造 30+ 条「微改写」/「重命名甲乙方」/
    「调整条款顺序」的合成样本，回归确保 \texttt{counts} 的稳定性。
  \item \textbf{Skill 元数据校验}：把 \texttt{install\_into\_openclaw.py} 加上
    YAML 解析 + \texttt{description} 长度 / 触发词覆盖度的预检，避免装入
    破损的 SKILL.md。
\end{enumerate}

% =============================================================
\section*{附录 A：版本要点}

\begin{itemize}[leftmargin=1.4em,itemsep=1pt]
  \item \textbf{v1.0}（2026-04）：单 skill；范本匹配器单门槛 $\theta=0.40$。
  \item \textbf{v1.1}（2026-05）：双 skill 拆分；匹配器升级为
    形状 $+$ 锚点双门槛；外移 QCC / ArkClaw 长篇手册；新增 CHANGELOG。
\end{itemize}

\end{document}
